% ID: SEL-D0024-OA-E04
% ESTADO: cotejada. Segunda pasada visual completa y borrador compilado revisados.
% Fuente: Selectividad/matematicas_aplicadas/reserva_b_EXAMEN 2-ANDALUCÍA.pdf
% SHA-256: 98b7723232acd3143de28397f7641923ed48292db7c35883ebd2b56a72b30b79
% Página: 1; opción A; ejercicio original 4.

En un centro docente hay 160 alumnos matriculados en 1\textordmasculine{} de ESO, 120 en 2\textordmasculine{}, 120 en 3\textordmasculine{}, 80 en 4\textordmasculine{}, 240 en 1\textordmasculine{} de Bachillerato y 200 en 2\textordmasculine{}. Se quiere constituir una comisión en la que todos los cursos estén representados de forma proporcional.
\begin{enumerate}[label=\alph*)]
\item \textbf{(1.25 puntos)} \textquestiondown{}Cuántos alumnos debe haber en la comisión y cuántos de cada curso si dicha comisión está formada por el 5\% del total del alumnado?
\item \textbf{(1.25 puntos)} \textquestiondown{}Cuál sería la composición de la comisión si queremos que haya 9 alumnos de 2\textordmasculine{} de ESO?
\end{enumerate}
